% titel: Blatt 1, A8: Molmasse aus Isotopenverteilung / A9: Größte Atomzahl
% skript: Mortimer Kap. 2 (Isotopenhäufigkeit, Mol)

% Teil I: Werkzeuge, kein Lösungsweg.

Aus der Häufigkeitsverteilung eines Elements seine (mittlere) Molmasse
berechnen und Proben verschiedener Stoffe über die Stoffmenge $n$ vergleichen.

\begin{werkzeugkasten}[Mittlere Masse und Stoffmenge]
Mittlere Atommasse aus Isotopenhäufigkeiten:
\begin{equation}
  \bar M = \sum_i f_i\,M_i
\end{equation}
($f_i$ = relative Häufigkeit in %, $M_i$ = Isotopenmasse). Die Molmasse ist
damit die gewichtete Summe.

Stoffmenge: $n = m/M$ (Masse durch Molmasse). Vergleich der \emph{Anzahl
Atome} = Vergleich von $n\cdot N_A$; da $N_A$ konstant ist, reicht der
Vergleich von $n$.

\textit{Skript:} Mortimer Kap. 2 (Isotopenhäufigkeit) und Kap. 3 (Mol,
Molmasse), Übung 1/8, 1/9.
\end{werkzeugkasten}

\begin{achtung}[Darauf achten]
\begin{itemize}
  \item Häufigkeiten als \emph{Bruch} (z. B.\ 0.917) einsetzen, nicht als Prozent.
  \item Molmasse $M$ je Stoff: für Wasser $M(\mathrm{H_2O})=18\ \mathrm{g/mol}$,
        Helium $4\ \mathrm{g/mol}$, Benzol $\mathrm{C_6H_6}$ $78\ \mathrm{g/mol}$,
        Eisen $55.8\ \mathrm{g/mol}$.
  \item Die Probe mit dem größten $n$ (nicht der größten Masse) hat die
        meisten Atome.
  \item Die Isotopenanteile in A8 aus der Grafik ablesen, nicht erraten.
\end{itemize}
\end{achtung}
